
\documentclass{article}
\usepackage{colortbl}
\usepackage{makecell}
\usepackage{multirow}
\usepackage{supertabular}

\begin{document}

\newcounter{utterance}

\centering \large Interaction Transcript for game `cladder', experiment `full\_v1.5\_default', episode 1290 with refineonly{-}Qwen3.5{-}27B.
\vspace{24pt}

{ \footnotesize  \setcounter{utterance}{1}
\setlength{\tabcolsep}{0pt}
\begin{supertabular}{c@{$\;$}|p{.15\linewidth}@{}p{.15\linewidth}p{.15\linewidth}p{.15\linewidth}p{.15\linewidth}p{.15\linewidth}}
    \# & \multicolumn{2}{c}{Player} && \multicolumn{2}{c}{Game Master} \\
    \hline

    \theutterance \stepcounter{utterance}  
    & & & \multicolumn{4}{p{0.6\linewidth}}{
        \cellcolor[rgb]{0.9,0.9,0.9}{
            \makecell[{{p{\linewidth}}}]{
                \texttt{\tiny{[P1$\langle$GM]}}
                \texttt{You are an expert in causal inference. The following question is not a typical commonsense query, but rather a meticulously designed question created by a professor specializing in causal inference, intended to assess the students' mastery of the course content.} \\
\texttt{Imagine a self{-}contained, hypothetical world with only the following conditions, and without any unmentioned factors or causal relationships: Gender has a direct effect on smoking and lung cancer. Smoking has a direct effect on tar deposit. Tar deposit has a direct effect on lung cancer. Gender is unobserved. For nonsmokers, the probability of high tar deposit is 36\%. For smokers, the probability of high tar deposit is 84\%. For nonsmokers and with no tar deposit, the probability of lung cancer is 12\%. For nonsmokers and with high tar deposit, the probability of lung cancer is 44\%. For smokers and with no tar deposit, the probability of lung cancer is 15\%. For smokers and with high tar deposit, the probability of lung cancer is 50\%. The overall probability of smoking is 67\%. Will smoking increase the chance of lung cancer?} \\
\texttt{Start your answer with 'Yes' or 'No', followed by additional reasoning or evidence to support your explanation.} \\
            }
        }
    }
     \\ \\

    \theutterance \stepcounter{utterance}  
    & \multicolumn{4}{p{0.6\linewidth}}{
        \cellcolor[rgb]{0.561,0.737,0.561}{
            \makecell[{{p{\linewidth}}}]{
                \texttt{\tiny{[P1$\rangle$GM]}}
                \texttt{Yes} \\
\\ 
\texttt{To determine whether smoking increases the chance of lung cancer, we must calculate the average causal effect of smoking on lung cancer, defined as \$P(Y=1 | do(S=1)) {-} P(Y=1 | do(S=0))\$, where \$Y\$ is lung cancer and \$S\$ is smoking.} \\
\\ 
\texttt{In this hypothetical world, the causal structure is:} \\
\texttt{1.  \$G \textbackslash{}to S\$ (Gender affects Smoking)} \\
\texttt{2.  \$G \textbackslash{}to Y\$ (Gender affects Lung Cancer)} \\
\texttt{3.  \$S \textbackslash{}to T\$ (Smoking affects Tar Deposit)} \\
\texttt{4.  \$T \textbackslash{}to Y\$ (Tar Deposit affects Lung Cancer)} \\
\\ 
\texttt{Here, \$G\$ (Gender) is an unobserved confounder between \$S\$ and \$Y\$. However, \$T\$ (Tar Deposit) is a mediator. The problem provides the conditional probabilities \$P(Y | S, T)\$ and \$P(T | S)\$. Since \$T\$ is a mediator that captures the effect of \$S\$ on \$Y\$ (and the provided \$P(Y|S,T)\$ accounts for the confounding effect of \$G\$ within each stratum of \$S\$ and \$T\$), we can calculate the interventional probability using the Law of Total Probability over the mediator \$T\$.} \\
\\ 
\texttt{The formula for the probability of lung cancer under the intervention \$do(S=s)\$ is:} \\
\texttt{\$\$P(Y=1 | do(S=s)) = \textbackslash{}sum\_\{t \textbackslash{}in \textbackslash{}\{0, 1\textbackslash{}\}\} P(Y=1 | S=s, T=t) \textbackslash{}cdot P(T=t | S=s)\$\$} \\
\\ 
\texttt{**Step 1: Calculate \$P(Y=1 | do(S=1))\$ (Probability of lung cancer if everyone smokes)**} \\
\texttt{We are given:} \\
\texttt{{-}   \$P(T=1 | S=1) = 0.84\$ (Probability of high tar for smokers)} \\
\texttt{{-}   \$P(T=0 | S=1) = 1 {-} 0.84 = 0.16\$} \\
\texttt{{-}   \$P(Y=1 | S=1, T=1) = 0.50\$} \\
\texttt{{-}   \$P(Y=1 | S=1, T=0) = 0.15\$} \\
\\ 
\texttt{\$\$P(Y=1 | do(S=1)) = (0.50 \textbackslash{}times 0.84) + (0.15 \textbackslash{}times 0.16)\$\$} \\
\texttt{\$\$P(Y=1 | do(S=1)) = 0.42 + 0.024 = 0.444\$\$} \\
\\ 
\texttt{**Step 2: Calculate \$P(Y=1 | do(S=0))\$ (Probability of lung cancer if no one smokes)**} \\
\texttt{We are given:} \\
\texttt{{-}   \$P(T=1 | S=0) =} \\
            }
        }
    }
    & & \\ \\

    \theutterance \stepcounter{utterance}  
    & & & \multicolumn{2}{p{0.3\linewidth}}{
        \cellcolor[rgb]{0.9,0.9,0.9}{
            \makecell[{{p{\linewidth}}}]{
                \texttt{\tiny{[GM$|$GM]}}
                \texttt{Yes} \\
            }
        }
    }
    & & \\ \\

    \theutterance \stepcounter{utterance}  
    & & & \multicolumn{2}{p{0.3\linewidth}}{
        \cellcolor[rgb]{0.9,0.9,0.9}{
            \makecell[{{p{\linewidth}}}]{
                \texttt{\tiny{[GM$|$GM]}}
                \texttt{yes} \\
            }
        }
    }
    & & \\ \\

    \theutterance \stepcounter{utterance}  
    & & & \multicolumn{2}{p{0.3\linewidth}}{
        \cellcolor[rgb]{0.9,0.9,0.9}{
            \makecell[{{p{\linewidth}}}]{
                \texttt{\tiny{[GM$|$GM]}}
                \texttt{game\_result = WIN} \\
            }
        }
    }
    & & \\ \\

\end{supertabular}
}

\end{document}
